%%%%%%%%%%%%%%%%%%%%%%%%%%%%%%%%%%%%%%%%%%%%%%%%%%%%%%%%%%%%%%%%%%%%%%
\section{Feuille de route prévisionnelle}
%%%%%%%%%%%%%%%%%%%%%%%%%%%%%%%%%%%%%%%%%%%%%%%%%%%%%%%%%%%%%%%%%%%%%%

Le projet ALORA est envisagé comme un projet de recherche progressif.
L'objectif n'est pas de développer immédiatement un système complet, mais
de construire, étape par étape, les fondations nécessaires au
développement d'outils fiables pour les langues orales peu dotées.

La feuille de route présentée ci-dessous constitue une première
planification susceptible d'évoluer en fonction des résultats obtenus et
des recommandations du laboratoire.

%%%%%%%%%%%%%%%%%%%%%%%%%%%%%%%%%%%%%%%%%%%%%%%%%%%%%%%%%%%%%%%%%%%%%%

\subsection{Phase 1 : Été 2026}

Cette première phase est consacrée à la construction des fondations du
projet.

Les principaux objectifs sont les suivants :

\begin{itemize}

\item réaliser une étude bibliographique sur les langues peu dotées et
les méthodes modernes de traitement de la parole ;

\item concevoir l'architecture générale du projet ALORA ;

\item développer une première plateforme de collecte collaborative ;

\item définir la structure du corpus et les métadonnées associées ;

\item mettre en place un pipeline de validation et de préparation des
données ;

\item constituer une première version du corpus oral.

\end{itemize}

Les livrables attendus sont :

\begin{itemize}

\item une plateforme fonctionnelle ;

\item un premier corpus structuré ;

\item une documentation technique ;

\item une méthodologie de collecte reproductible.

\end{itemize}

%%%%%%%%%%%%%%%%%%%%%%%%%%%%%%%%%%%%%%%%%%%%%%%%%%%%%%%%%%%%%%%%%%%%%%

\subsection{Phase 2 : Semestre S5}

Cette phase sera consacrée aux premières expérimentations en
apprentissage automatique.

Les travaux envisagés comprennent notamment :

\begin{itemize}

\item l'analyse exploratoire du corpus ;

\item le prétraitement avancé des données audio ;

\item l'étude de modèles de représentation de la parole ;

\item les premières expérimentations sur la découverte automatique de
structures phonologiques ;

\item l'évaluation des approches retenues.

\end{itemize}

L'objectif est de disposer, à la fin du semestre, d'une première base
expérimentale permettant d'orienter les développements futurs.

%%%%%%%%%%%%%%%%%%%%%%%%%%%%%%%%%%%%%%%%%%%%%%%%%%%%%%%%%%%%%%%%%%%%%%

\subsection{Phase 3 : Projet de Fin d'Études}

Le Projet de Fin d'Études constituera la phase principale de
développement scientifique.

Selon les résultats obtenus durant les étapes précédentes, plusieurs
axes pourront être approfondis :

\begin{itemize}

\item amélioration des représentations de la parole ;

\item découverte automatique des unités phonologiques ;

\item apprentissage progressif de la structure de la langue ;

\item étude de méthodes d'induction d'un système d'écriture ;

\item développement de premiers démonstrateurs (transcription,
recherche vocale, traduction ou autres applications).

\end{itemize}

Cette phase permettra également de consolider les résultats obtenus dans
une perspective de publication ou de poursuite en doctorat.

%%%%%%%%%%%%%%%%%%%%%%%%%%%%%%%%%%%%%%%%%%%%%%%%%%%%%%%%%%%%%%%%%%%%%%

\subsection{Perspective à long terme}

Au-delà du cas du soussou, l'ambition d'ALORA est de proposer une
méthodologie générale pouvant être adaptée à d'autres langues orales peu
dotées.

À terme, cette approche pourrait contribuer à la création de nouvelles
ressources linguistiques et faciliter le développement d'outils
numériques pour des langues aujourd'hui très peu représentées dans les
technologies de traitement automatique de la parole.

%%%%%%%%%%%%%%%%%%%%%%%%%%%%%%%%%%%%%%%%%%%%%%%%%%%%%%%%%%%%%%%%%%%%%%

\begin{objectif}

Cette feuille de route ne constitue pas un planning figé.

Elle représente une première proposition de travail destinée à servir de
base de discussion avec le laboratoire. Les différentes étapes pourront
être adaptées en fonction des recommandations scientifiques, des
contraintes rencontrées et des résultats obtenus au cours du projet.

\end{objectif}

%%%%%%%%%%%%%%%%%%%%%%%%%%%%%%%%%%%%%%%%%%%%%%%%%%%%%%%%%%%%%%%%%%%%%%

\section*{Conclusion}

\addcontentsline{toc}{section}{Conclusion}

Le projet ALORA est né d'une réflexion sur les difficultés rencontrées
par les langues orales peu dotées dans le contexte actuel de
l'intelligence artificielle.

Cette note de cadrage présente une première vision du projet ainsi que
les travaux envisagés pour sa phase initiale.

Je suis conscient que plusieurs choix méthodologiques restent ouverts et
que certaines hypothèses pourront évoluer au cours des expérimentations.
C'est précisément pour cette raison que je souhaite bénéficier d'un
encadrement scientifique dès les premières étapes du projet.

J'espère que cette proposition pourra constituer le point de départ
d'une collaboration enrichissante et je serais très heureux d'échanger
avec vous afin de recueillir vos remarques et vos conseils pour faire
évoluer cette idée vers un véritable projet de recherche.